\exercisesection{\S~8}

\begin{enumerate}
\def\labelenumi{\arabic{enumi})}
\item
  令 \(f\) 为 \(\mathfrak{g}\) 上的不变多项式函数。证明 \(f\) 在 \(\operatorname{Aut}(\mathfrak{g})\) 下不变，当且仅当 \(f|\mathfrak{h}\) 在 \(\operatorname{Aut}(\mathrm{R})\) 下不变。推出，若 \(\mathrm{R}\) 的邓金图有非平凡自同构，则 \(\mathfrak{g}\) 上存在一个在 \(\operatorname{Aut}(\mathfrak{g})\) 下不变的不变多项式函数。
\item
  取 \(\mathfrak{g} = \mathfrak{sl}(3,k)\)。证明 \(x\mapsto \det (x)\) 是 \(\mathfrak{g}\) 上的不变多项式函数，但在 \(\operatorname {Aut}(\mathfrak{g})\) 下不变（使用自同构 \(x\mapsto -^t x\)）。
\item
  令 \(\mathfrak{a}\) 为半单李代数，且 \(s\in \mathrm{Aut}(\mathfrak{a})\)。证明以下条件等价：
\end{enumerate}

\begin{enumerate}
\def\labelenumi{(\roman{enumi})}
\item
  \(s \in \text{Aut}_{0}(\mathfrak{a})\)。
\item
  \(s\) 在 \(\mathrm{U}(\mathfrak{a})\) 的中心上平凡作用。
\item
  对所有 \(x \in \mathfrak{a}\)，存在 \(t \in \mathrm{Aut}_0(\mathfrak{a})\)，使得 \(tx = sx\)。
\end{enumerate}

（使用命题6证明 (iii) \(\Longrightarrow\) (i)。）

\begin{enumerate}
\def\labelenumi{\arabic{enumi})}
\setcounter{enumi}{3}
\item
  证明在命题2的推论2和定理1 (ii) 中，可以只考虑权为根权的表示 \(\rho\)（注意，当 \(k[\mathrm{P}]^{\mathrm{W}}\) 被 \(k[\mathrm{Q}]^{\mathrm{W}}\) 替代时，命题1仍然成立，其中 \(\mathrm{Q}\) 是根权群）。
\item
  沿用§6和§7的记号。令 \(\lambda \in \mathfrak{h}^*\)。若对任意 \(w \in \mathrm{W}\) 且 \(w \neq 1\)，都有 \((\lambda + \rho) - w(\lambda + \rho) \notin \mathrm{Q}_+\)，则 \(\mathrm{Z}(\lambda)\) 是单的。
\end{enumerate}

（使用定理2的推论1 (ii)。）

\begin{enumerate}
\def\labelenumi{\arabic{enumi})}
\setcounter{enumi}{5}
\item
  令 \(\mathfrak{a}\) 为李代数，\(f\) 为 \(\mathfrak{a}\) 上的多项式函数，\(x, y\) 为 \(\mathfrak{a}\) 的两个元素。置 \(f_y = \theta^*(y)f\)（参见~第3节），并以 \(\mathrm{D}_x f\) 表示 \(f\) 在 \(x\) 处的切线性映射（第VII章附录I第2节）。证明 \(f_y(x) = (\mathrm{D}_xf)([x,y])\)。推出，当 \(\mathrm{D}_xf\) 不变时，\(\operatorname{Im}(\operatorname{ad}x)\) 在 \(f\) 上为零。
\item
  *令 \(d_1, \ldots, d_l\) 为代数 I(g) 的特征次数，参见~第V章§5第1节。对任意整数 \(n \geq 0\)，以 \(r_n\) 表示 I(g) 上的 S(g) 的齐次基中次数为 \(n\) 的元素个数，并置 \(r(\mathrm{T}) = \sum_{n=0}^{\infty} r_n \mathrm{T}^n\)。证明
\end{enumerate}

\[
r (\mathrm{T}) = (1 - \mathrm{T}) ^ {- \mathrm{N}} \prod_ {i = 1} ^ {l} (1 - \mathrm{T} ^ {d _ {i}}), \quad \text { where } \mathrm{N} = \dim (\mathfrak {g}).
\]

\begin{enumerate}
\def\labelenumi{\arabic{enumi})}
\setcounter{enumi}{7}
\tightlist
\item
  置 \(l = \operatorname{rk}(\mathfrak{g}), N = \dim(\mathfrak{g})\)。若 \(x \in \mathfrak{g}\)，按下式定义 \(a_i(x)\)，其中 \(0 \leq i \leq N\)：
\end{enumerate}

\[
\det (\mathrm{T} + \operatorname{ad} x) = \sum_ {i = 0} ^ {\mathrm{N}} \mathrm{T} ^ {\mathrm{N} - i} a _ {i} (x).
\]

按上述方式定义的函数 \(a_{i}\) 是次数为 i 的齐次多项式，并且在 \(\operatorname{Aut}(\mathfrak{g})\) 下不变。若 \(x \in h\) 且 \(i \leq N - l\)，则 \(a_{i}(x)\) 是 \(\alpha(x)\)（\(\alpha \in R\)）的第 i 个初等对称函数；特别地，\(a_{N-l}(x) = \prod_{\alpha \in \mathbb{R}} \alpha(x)\)。构造一个例子，使得这些 \(a_{i}\) 不生成在 \(\operatorname{Aut}(\mathfrak{g})\) 下不变的 g 上的多项式函数代数。

\begin{enumerate}
\def\labelenumi{\arabic{enumi})}
\setcounter{enumi}{8}
\tightlist
\item
  记号同第5节。令 \(\lambda \in \mathfrak{h}^*\)，\(z \in \mathbf{Z}\)，\(z'\) 为 z 在 \(z\)\(\mathrm{U}(\mathfrak{g})\) 的主反自同构下的像，令 \(w_0\) 为将 B 变换为 \(-\mathrm{B}\) 的 W 中元素。证明 \(\chi_{\lambda}(z) = \chi_{-w_0\lambda}(z')\)。
\end{enumerate}

（只需对 \(\lambda \in \mathrm{P}_{++}\) 证明此事。考虑 z 在 \(z\)\(\mathrm{E}(\lambda)\) 和 \(\mathrm{E}(\lambda)^*\) 上的作用；使用§7的命题11。）

\begin{enumerate}
\def\labelenumi{\arabic{enumi})}
\setcounter{enumi}{9}
\tightlist
\item
  （本题及以下三题沿用§6的记号。）令 \(\lambda \in \mathfrak{h}^*\)。
\end{enumerate}

\begin{enumerate}
\def\labelenumi{\alph{enumi})}
\tightlist
\item
  令 \(\mathrm{N}, \mathrm{N}'\) 为 \(\mathfrak{g}\) 的 \(\mathrm{Z}(\lambda)\)-子模，满足 \(\mathrm{N}' \subset \mathrm{N}\)，且 \(\mathrm{N} / \mathrm{N}'\) 是单的。证明存在 \(\mu \in \lambda - \mathrm{Q}_+\)，使得 \(\mathrm{N} / \mathrm{N}'\) 同构于 \(\mathrm{E}(\mu)\)（应用§6的定理1），并且 \(\mu + \rho \in \mathrm{W}.(\lambda + \rho)\)（应用定理2的推论1）。
\item
  证明 \(\mathrm{Z}(\lambda)\) 有一个 Jordan-Hölder 序列。（应用 a)，以及 \(\mathrm{Z}(\lambda)\) 的权具有有限重数这一事实。）
\end{enumerate}

\begin{enumerate}
\def\labelenumi{\arabic{enumi})}
\setcounter{enumi}{10}
\tightlist
\item
  令 \(\lambda \in \mathfrak{h}^*\)，令 \(V\) 为 \(\mathfrak{g}\) 的子模且包含于 \(Z(\lambda)\)。证明存在 \(\mu \in \mathfrak{h}^*\)，使得 \(V\) 包含一个简单的 \(\mathfrak{g}\)-子模，它同构于 \(Z(\mu)\)。
\end{enumerate}

（令 A 为满足 V 包含一个同构于 \(\nu \in \mathfrak{h}^*\) 的 \(\mathfrak{g}\)-子模的 \(Z(\nu)\) 的集合。使用第6节的命题6，先证明 \(A \neq \varnothing\)。然后证明 A 是有限的，并取 A 中满足 \(\mu\) 的元素 \((\mu - Q_{+}) \cap A = \{\mu\}\)。）

¶ 12) a) 令 \(\lambda, \mu \in \mathfrak{h}^*\)。从 \(\mathfrak{g}\) 到 \(Z(\mu)\) 的每个非零 \(Z(\lambda)\)-同态都是单射。（使用第6节的命题6。）

\begin{enumerate}
\def\labelenumi{\alph{enumi})}
\setcounter{enumi}{1}
\item
  令 \(r \in \mathbf{N}\)，令 A 为 \(\mathbf{N}^r\) 的有限子集，且 \(m = \mathrm{Card}(\mathrm{A})\)。对于 \(\xi \in \mathbf{N}^r\)，令 \(s(\xi)\) 表示 \(\xi\) 的坐标之和，并令 \(\mathfrak{P}_{\mathrm{A}}(\xi)\) 表示满足 \((n_{\alpha})_{\alpha \in \mathrm{A}}\) 中每个整数 \(\geq 0\) 且 \(\xi = \sum_{\alpha \in \mathrm{A}} n_{\alpha}\alpha\) 的这类族的数目。则有 \(\mathfrak{P}_{\mathrm{A}}(\xi) \leq (s(\xi) + 1)^m\)，对所有 \(\xi \in \mathbf{N}^r\) 成立。（对 \(m\) 作归纳论证。）
\item
  令 \(\lambda, \mu \in \mathfrak{h}^*\)。证明 \(\dim \operatorname{Hom}_{\mathfrak{g}}(\mathrm{Z}(\mu), \mathrm{Z}(\lambda)) \leq 1\)。（令 \(\varphi_1\) 和 \(\varphi_2\) 为非零 \(\mathfrak{g}\)-同态，从 \(\mathrm{Z}(\mu)\) 到 \(\mathrm{Z}(\lambda)\)。若 \(\operatorname{Im}(\varphi_1) = \operatorname{Im}(\varphi_2)\)，则由§6的命题5，\(\varphi_1\) 和 \(\varphi_2\) 线性相关。设 \(\operatorname{Im}(\varphi_1) \neq \operatorname{Im}(\varphi_2)\)。若 \(\mathrm{Z}(\mu)\) 是单的，则 \(\operatorname{Im}(\varphi_1) + \operatorname{Im}(\varphi_2)\) 是直和；推出 \(\mathfrak{P}(\xi + \lambda - \mu) \geq 2\mathfrak{P}(\xi)\) 对所有 \(\xi \in \mathfrak{h}^*\) 都成立，这与 \(b\) 矛盾。在一般情形下，使用相关习题。）
\end{enumerate}

当 \(\dim\operatorname{Hom}_{\mathfrak{g}}(\mathrm{Z}(\mu),\mathrm{Z}(\lambda))=1\) 时，滥用记号写作 \(\mathrm{Z}(\mu)\subset\mathrm{Z}(\lambda)\)。d) 令 \(\nu\in\mathfrak{h}^{*}\)。使得 \(\lambda\in\mathfrak{h}^{*}\) 的 \(\mathrm{Z}(\lambda-\nu)\subset\mathrm{Z}(\lambda)\) 的集合在 \(\mathfrak{h}^{*}\) 的 Zariski 拓扑中是闭集。

¶ 13) a) 令 \(\mathfrak{a}\) 为幂零李代数，\(x \in \mathfrak{a}, n \in \mathbf{N}, p \in \mathbf{N}\)。存在 \(l \in \mathbf{N}\)，使得对所有 \(x^{l}y_{1} \ldots y_{n} \in \mathrm{U}(\mathfrak{a})x^{p}\)，都有 \(y_{1}, \ldots, y_{n} \in \mathfrak{a}\)。

\begin{enumerate}
\def\labelenumi{\alph{enumi})}
\setcounter{enumi}{1}
\tightlist
\item
  令 \(\lambda, \mu \in \mathfrak{h}^*\)，且 \(\alpha \in \mathrm{B}\) 满足
\end{enumerate}

\[
\mathrm{Z} (s _ {\alpha} \mu - \rho) \subset \mathrm{Z} (\mu - \rho) \subset \mathrm{Z} (\lambda - \rho).
\]

假设 \(\lambda \in \mathrm{P}\)。令 \(p = \lambda (H_{\alpha})\in \mathbf{Z}\)。证明：

\(b_{1})\) 若 \(p \leq 0\)，则 \(\mathrm{Z}(\lambda - \rho) \subset \mathrm{Z}(s_{\alpha}\lambda - \rho)\)。

\[
b _ {2}) \text {   If   } p > 0, \text {   then   } \mathrm{Z} (s _ {\alpha} \mu - \rho) \subset \mathrm{Z} (s _ {\alpha} \lambda - \rho) \subset \mathrm{Z} (\lambda - \rho).
\]

（使用 \(a\)，以及§6命题6的推论1。）

\begin{enumerate}
\def\labelenumi{\alph{enumi})}
\setcounter{enumi}{2}
\tightlist
\item
  令 \(\lambda \in \mathfrak{h}^*\)，\(\alpha \in \mathrm{R}\)，且 \(m = \lambda(H_{\alpha})\)。假设 \(m \in \mathbf{N}\)。证明
\end{enumerate}

\[
\mathrm{Z} (s _ {\alpha} \lambda - \rho) \subset \mathrm{Z} (\lambda - \rho).
\]

(先对 \(\lambda \in \mathrm{P}\) 使用 \(b\) 证明此事，然后在一般情形下使用习题12 \(d\)。) \(^{17}\)

\begin{enumerate}
\def\labelenumi{\arabic{enumi})}
\setcounter{enumi}{13}
\item
  *令 \(\mathfrak{a}\) 为半单李代数，令 \(Z(\mathfrak{a})\) 为 \(U(\mathfrak{a})\) 的中心。证明 \(U(\mathfrak{a})\) 是自由的 \(Z(\mathfrak{a})\)-模。（注意 \(grU(\mathfrak{a})\) 同构于 \(S(\mathfrak{a})\) 和 \(S(\mathfrak{a}^*)\)，并使用第3节的注2。）*
\item
  令 \(x\) 为 \(\mathfrak{g}\) 的一个可对角化元素（§3，习题10），令 \(y\) 为 \(\mathfrak{g}\) 的一个半单元素，满足对 \(f(x) = f(y)\) 上的每个不变多项式函数 \(f\) 都有 \(\mathfrak{g}\)。证明存在 \(s \in \mathrm{Aut}_e(\mathfrak{g})\)，使得 \(sy = x\)。（注意，ad \(x\) 和 ad \(y\) 有相同的特征多项式，参见~习题8，因而 y 是可对角化的。接着使用§3的习题10，将情形化归为 x 和 y 包含于 h 中，再使用定理1 (i) 和引理6证明 x 和 y 在 W 下共轭。）
\item
  令 \(\mathfrak{a}\) 为半单李代数，\(x\) 为 \(\mathfrak{a}\) 的一个元素，\(x_{s}\) 为 \(x\) 的半单部分。证明，若 \(f\) 是 \(\mathfrak{a}\) 上的不变多项式函数，则 \(f(x) = f(x_{s})\)。（将情形化归为 \(f\) 形如 \(x \mapsto \operatorname{Tr} \rho(x)^n\) 的情形。）
\item
  假设 \(k\) 是代数闭域，并置 \(\mathrm{G} = \operatorname{Aut}_e(\mathfrak{g})\)。令 \(x, y \in \mathfrak{g}\)。证明以下条件等价：
\end{enumerate}

\begin{enumerate}
\def\labelenumi{(\roman{enumi})}
\item
  \(x\) 和 \(y\) 的半单部分是 G-共轭的。
\item
  对 \(f\) 上的每个不变多项式函数 \(\mathfrak{g}\)，都有 \(f(x) = f(y)\)。（使用习题15和16。）
\end{enumerate}

\begin{enumerate}
\def\labelenumi{\arabic{enumi})}
\setcounter{enumi}{17}
\tightlist
\item
  令 \(\mathfrak{a}\) 为半单李代数，\(l = \mathrm{rk}(\mathfrak{a})\)，I 为 \(\mathfrak{a}\) 上不变多项式函数的代数，且 \(\mathrm{P}_1, \ldots, \mathrm{P}_l\) 为 I 的齐次元素，并且作为代数生成 I。\(\mathrm{P}_i\) 定义多项式映射 \(\mathrm{P} : \mathfrak{a} \to k^l\)。若 \(x \in \mathfrak{a}\)，以 \(\mathrm{D}_x\mathrm{P} : \mathfrak{a} \to k^l\) 表示 P 在 \(x\) 处的切线性映射（第VII章附录I第2节）。
\end{enumerate}

\begin{enumerate}
\def\labelenumi{\alph{enumi})}
\tightlist
\item
  令 \(\mathfrak{h}\) 为 \(\mathfrak{a}\) 的 Cartan 子代数，令 \(x \in \mathfrak{h}\)。证明以下条件等价：
\end{enumerate}

\begin{enumerate}
\def\labelenumi{(\roman{enumi})}
\item
  \(D_{x}P|\mathfrak{h}\) 是从 \(\mathfrak{h}\) 到 \(k^{l}\) 的同构；
\item
  \(x\) 是正则的。
\end{enumerate}

（将情形化归为分裂情形。取 \(\mathfrak{h}\) 的一个基，并记 \(d(x)\) 关于此基的矩阵的行列式为 \(\mathrm{D}_x\mathrm{P}|\mathfrak{h}\)。利用第V章§5第4节的命题5证明，存在 \(c \in k^*\)，使得 \(d(x)^2 = c \prod_{\alpha \in \mathrm{R}} \alpha(x)\)，

其中 \(\alpha\) 属于 \(R\) 的根集 \((\mathfrak{a},\mathfrak{h})\)。）

若满足这些条件，则 \(\mathfrak{a} = \mathfrak{h} \oplus \operatorname{Im} \operatorname{ad}(x)\)，且 \(\operatorname{Ker} D_x P = \operatorname{Im} \operatorname{ad}(x)\)（使用习题6证明 \(D_x P\) 在 \(\operatorname{Im} \operatorname{ad}(x)\) 上为零）。

\begin{enumerate}
\def\labelenumi{\alph{enumi})}
\setcounter{enumi}{1}
\tightlist
\item
  证明满足 \(x \in a\) 秩为 l 的 \(D_{x}P\) 的集合是 a 在 Zariski 拓扑中的稠密开子集。
\end{enumerate}

